% =============================================================================
% Section 5 — Conclusions
% =============================================================================
\section{Conclusions}
\label{sec:conclusions}

We have extended the fluid-filled viscoelastic oblate spheroidal shell
model of Mace and Mace~\cite{Mace2026browntone} to the sub-bass band
(\SIrange{20}{80}{\hertz}) to assess whether airborne acoustic coupling
at concert sound pressure levels can produce physiologically meaningful
tissue displacement in the human abdomen.  The principal findings are:
%
\begin{enumerate}
  \item The Helmholtz number $ka$ increases from $0.011$ at
    \SI{3.95}{\hertz} to $0.23$ at \SI{80}{\hertz}, improving the
    acoustic coupling factor $(ka)^2$ by up to ${\sim}400\times$.

  \item This improvement is entirely offset by the off-resonance
    attenuation of the SDOF transfer function: at \SI{40}{\hertz}
    ($r \approx 10$), $H \approx 0.01$.

  \item At peak EDM levels (\SI{115}{\decibel} at \SI{40}{\hertz}),
    airborne-induced tissue displacement is approximately four orders
    of magnitude below whole-body vibration perception thresholds.

  \item No ``coupling transition band'' exists in the
    \SIrange{1}{100}{\hertz} range at concert SPLs: the displacement
    never exceeds \SI{1}{\percent} of the perception threshold at any
    frequency.

  \item The familiar sensation of ``feeling the bass'' at concerts is
    not attributable to airborne acoustic coupling to the abdominal
    cavity.  It must arise from structural vibration (floor, seats),
    chest-wall coupling, somatosensory stimulation, or a combination
    thereof.
\end{enumerate}

These results close an important gap between structural acoustics and
concert psychoacoustics.  While the improved coupling at sub-bass
frequencies might initially suggest a pathway for airborne visceral
excitation, the off-resonance penalty is decisive.  Sound engineers and
venue designers seeking to maximise the physical impact of bass should
focus on structural vibration transmission rather than acoustic levels.
